\section{背景、工作负载与设计动机}
\label{sec:background}

\subsection{卷积循环与 KCS 分块}

对输入张量 $X\in\mathbb{Z}^{H\times W\times C_{in}}$、权重
$W\in\mathbb{Z}^{K_h\times K_w\times C_{in}\times C_{out}}$，输出为
\begin{equation}
Y[o_y,o_x,c_o]=b[c_o]+\sum_{k_y,k_x,c_i}
 X[i_y,i_x,c_i]W[k_y,k_x,c_i,c_o],
\end{equation}
其中 $i_y=o_y s+k_y-p$、$i_x=o_x s+k_x-p$。将
$K=K_hK_wC_{in}$，并把 $(o_y,o_x)$ 展平为空间索引 $S$，卷积可视作
$S\times K$ 与 $K\times C_{out}$ 的矩阵乘。本文将三个有限硬件维度记为
KCS：
\begin{itemize}[leftmargin=1.8em]
  \item $K$ 维按阵列行数 $R=18$ 分为 $N_K=\lceil K/R\rceil$ 个 pass；
  \item $C_{out}$ 按 $C_T=32$ 分为 $N_C=\lceil C_{out}/C_T\rceil$ 个块；
  \item 空间 $S$ 按层自适应 tile 高度分为 $N_S$ 个 tile。
\end{itemize}
于是一次硬件 dispatch 对一个空间 tile 依次完成 $N_C N_K$ 个计算上下文。
第一 pass 从 bias 初始化 PSUM，后续 pass 读取上一次结果：
\begin{align}
P_0[s,c]&=b[c]+\sum_{k=0}^{R-1}X_s[k]W[k,c],\\
P_t[s,c]&=P_{t-1}[s,c]+\sum_{k=tR}^{(t+1)R-1}X_s[k]W[k,c].
\end{align}
只有 $t=N_K-1$ 的结果进入重量化和 HWC 写回。这样的分块允许固定阵列
覆盖不同 $K$ 和 $C_{out}$，代价是必须高效、严格地处理输入重放和 PSUM
反馈。

\subsection{完整双尺度工作负载}

本文采用 Ultralytics v9.5.0 对应的 COCO80 YOLOv3-tiny 图结构
\cite{ultralyticsv95,coco2014}。检测头每个 anchor 输出
$4+1+80=85$ 个值，每个尺度有 3 个 anchor，因而 P4/P5 均为 255 通道。
P4 作用于 26$\times$26、stride 16 特征，P5 作用于 13$\times$13、stride 32
特征，两头共 $26^2\times3+13^2\times3=2535$ 个 anchor。完整卷积层配置
由生成器从 \code{tools/coco80/model\_spec.json} 读取，见
\cref{tab:layer-config}。

\begin{table*}[t]
\centering
\caption{完整 YOLOv3-tiny 的 13 个 PL 卷积层及安全分块。$K_{\mathrm{pass}}=\lceil K/18\rceil$，输出块宽为 32。}
\label{tab:layer-config}
\small
\setlength{\tabcolsep}{4.2pt}
\begin{tabular}{lrrrrrrrr}
\toprule
层 & $H\times W$ & $C_{in}$ & $C_{out}$ & 核 & $K$ & $K$-pass & Cout块 & tile\_h/块数 \\
\midrule
m0 & 416$\times$416 & 3 & 16 & 3$\times$3 & 27 & 2 & 1 & 2 / 208 \\
m2 & 208$\times$208 & 16 & 32 & 3$\times$3 & 144 & 8 & 1 & 4 / 52 \\
m4 & 104$\times$104 & 32 & 64 & 3$\times$3 & 288 & 16 & 2 & 8 / 13 \\
m6 & 52$\times$52 & 64 & 128 & 3$\times$3 & 576 & 32 & 4 & 8 / 7 \\
m8 & 26$\times$26 & 128 & 256 & 3$\times$3 & 1152 & 64 & 8 & 13 / 2 \\
m10 & 13$\times$13 & 256 & 512 & 3$\times$3 & 2304 & 128 & 16 & 13 / 1 \\
m13 & 13$\times$13 & 512 & 1024 & 3$\times$3 & 4608 & 256 & 32 & 8 / 2 \\
m14 & 13$\times$13 & 1024 & 256 & 1$\times$1 & 1024 & 57 & 8 & 13 / 1 \\
m15 & 13$\times$13 & 256 & 512 & 3$\times$3 & 2304 & 128 & 16 & 13 / 1 \\
m16 & 13$\times$13 & 256 & 128 & 1$\times$1 & 256 & 15 & 4 & 13 / 1 \\
m19 & 26$\times$26 & 384 & 256 & 3$\times$3 & 3456 & 192 & 8 & 6 / 5 \\
p4\_detect & 26$\times$26 & 256 & 255 & 1$\times$1 & 256 & 15 & 8 & 13 / 2 \\
p5\_detect & 13$\times$13 & 512 & 255 & 1$\times$1 & 512 & 29 & 8 & 13 / 1 \\
\bottomrule
\end{tabular}
\end{table*}


层间差异并非简单的“越深越慢”。m0 的 $K=27$、$N_K=2$，但 416 宽输入
使空间 tile 和 HWC 搬运占主导；m13 的空间仅 13$\times$13，却有
$K=4608$ 和 256 个 pass；m19 由上采样分支与 route 分支拼接得到
$C_{in}=384$，$K=3456$，其 materialized entry 数为
$26\times6\times192=29952$。若将 m19 的 tile 高度从 6 增至 7，条目数变为
34944，超过 32768 深度，因而层自适应表不仅是性能启发式，也是容量安全
合同。两个检测头的 255 通道还要求输出打包器在最后块屏蔽 1 个 lane，并
以 TKEEP 保持跨像素和跨 tile 的连续字节流。

\subsection{HWC 布局为何成为阵列问题}

CPU、图像库和许多边缘推理软件偏好 HWC/NHWC，因为一个像素的通道连续，
便于预处理、池化、拼接和 DMA burst。矩阵乘阵列则希望同一空间位置的
$R$ 个 $K$ 元素同时到达不同行。最直接的方法是由 A53 对每个卷积生成
im2col，但这会写出 $S\times K$ 的展开张量；另一种方法是每个 pass 从 DDR
重新读取 HWC，并在 PL 即时生成对应 18 个 lane，这又使同一 tile 被读取
$N_K$ 次。对于 m13，二者都会把“输入一次”放大为与 256 pass 相关的布局
和存储开销。

相关研究已展示行缓存、数据流搜索和布局协同可以显著改善复用
\cite{fpgaConvNet2015,dnnweaver2018,snowflake2017,feather2024}。本文的设计点
在于把 HWC 语义保留到 PL 内部：DMA 输入仍是连续 HWC 字节，物化器按
输出坐标、卷积核位置和通道生成 18-lane 向量，并将当前 spatial tile 的
完整 $K$-pass 向量集写入可重放缓存。这样软件不需要知道阵列 bank 布局，
硬件也不需要在每个 pass 重新访问 DDR。

\subsection{串行时间线的空泡来源}

传统逐 pass 控制可表示为四个串行阶段：加载权重 $W_t$、填充输入
$I_t$、计算 $C_t$、排出 PSUM $D_t$。总时间近似为
\begin{equation}
T_{legacy}=\sum_{t=0}^{N_K-1}(T_{W,t}+T_{I,t}+T_{C,t}+T_{D,t})+T_{switch}.
\end{equation}
当 $N_K$ 很大时，权重与反馈等待会被重复数百次；当空间 tile 较大时，
重复 IFM 填充同样昂贵。理想重叠并不是简单取四者最大值，因为 $D_t$ 的
结果是 $C_{t+1}$ 的数据依赖，而且权重、IFM、PSUM 缓冲可能被上下文共享。
要接近
\begin{equation}
T_{LASA}\approx T_{materialize}+\sum_t\max(T_{W,t+1},T_{C,t},T_{D,t})
\end{equation}
必须同时解决数据依赖和缓冲所有权。

因此，本文把调度优化拆为四个可验证机制：权重预载隐藏 $T_W$；HWC
materialize-once/replay-many 去除重复 $T_I$；continuous PSUM 与 drain ring
缩短 $D_t\rightarrow C_{t+1}$；tagged fast handoff 隐藏 $T_{switch}$。
\cref{fig:timeline} 在机制层面对比两条时间线。本文不会以该概念图替代受控
消融，消融数值必须由相同源码、阵列、频率和输入的独立 bitstream 得到。

\begin{figure*}[t]
\centering
\begin{tikzpicture}[x=0.8cm,y=0.8cm,font=\small]
\node[anchor=east] at (0,2.8) {逐 pass 基线};
\foreach \x/\w/\c/\t in {0/2/LasaGreen/IFM,2/1/LasaOrange/W,3/3/LasaBlue/C,6/1/LasaRed/D,7/2/LasaGreen/IFM,9/1/LasaOrange/W,10/3/LasaBlue/C,13/1/LasaRed/D}{
  \fill[\c!65] (\x,2.35) rectangle +(\w,0.7); \node at (\x+\w/2,2.7) {\t};
}
\draw[->] (0,2.2)--(14.5,2.2) node[right]{time};
\node[anchor=east] at (0,1.0) {\LASA};
\fill[LasaGreen!65] (0,0.55) rectangle +(2,0.7); \node at (1,0.9){物化};
\fill[LasaOrange!65] (1.2,0.55) rectangle +(1.2,0.7); \node at (1.8,0.9){W0};
\fill[LasaBlue!65] (2.4,0.55) rectangle +(3,0.7); \node at (3.9,0.9){C0};
\fill[LasaOrange!65] (4.2,1.28) rectangle +(1.2,0.45); \node at (4.8,1.5){W1};
\fill[LasaRed!65] (5.4,0.55) rectangle +(0.7,0.7); \node at (5.75,0.9){D};
\fill[LasaBlue!65] (6.1,0.55) rectangle +(3,0.7); \node at (7.6,0.9){C1/replay};
\fill[LasaOrange!65] (7.8,1.28) rectangle +(1.2,0.45); \node at (8.4,1.5){W2};
\fill[LasaRed!65] (9.1,0.55) rectangle +(0.7,0.7); \node at (9.45,0.9){D};
\fill[LasaBlue!65] (9.8,0.55) rectangle +(3,0.7); \node at (11.3,0.9){C2/replay};
\fill[LasaOrange!65] (11.3,1.28) rectangle +(1.5,0.45); \node at (12.05,1.5){next ctx};
\draw[->] (0,0.4)--(14.5,0.4) node[right]{time};
\draw[decorate,decoration={brace,amplitude=4pt},thick] (2,3.2)--(3,3.2) node[midway,above=4pt]{空泡};
\draw[decorate,decoration={brace,amplitude=4pt},thick] (7,3.2)--(10,3.2) node[midway,above=4pt]{重复搬运};
\node[anchor=west,align=left] at (0,-0.5) {W：权重加载；C：阵列计算；D：PSUM drain。上行小条表示预载与当前计算重叠。};
\end{tikzpicture}
\caption{逐 pass 搬运与 \LASA 跨 pass 调度的概念时间线。\LASA 只物化一次
IFM tile，后续 pass 由片上重放；权重预载、连续 PSUM 和下一上下文准备与当前
计算重叠。该图表达机制，不以未完成的消融数值替代测量。}
\label{fig:timeline}
\end{figure*}


\subsection{相关工作的能力维度}

现有工作大致可分为四类。固定数据流加速器通过 weight-stationary、
row-stationary 或 output-stationary 强化特定复用\cite{chen2016eyeriss,
jouppi2017tpu,sze2017efficient}；可重构数据流和编译器工作在层间改变循环
映射\cite{maeri2018,eyerissv2,simba2019,nvdla2018}；FPGA 框架强调定制精度、
HLS 生成和层级流水\cite{finn2017,fpgaConvNet2015,dnnweaver2018,
vitisai2022}；目标检测专用工作则把 YOLO 的卷积和后处理映射到一块或多块
FPGA\cite{multicoreyolo2023,tinyyoloaccel2022,yolov3fpga2019,
yolov2fpga2018}。

本文不以“峰值 GOPS 最大”作为差异。其能力组合是：面向外部 HWC 的片上
窗口物化；同一 tile 跨 $K$-pass 重放；带身份与所有权的上下文重叠；完整
双尺度网络中 13 个卷积的统一执行；以及由 RTL-semantic host、板端中间
节点、COCO 指标和 artifact hash 构成的端到端证据链。\cref{tab:related}
以能力而非孤立峰值总结相关方向。

固定数据流研究证明了空间阵列可以通过选择合适的驻留对象获得高数据复用，
但其分析通常假设进入加速器的数据已经满足内部布局。对于边缘 MPSoC，这一
假设会把代价转移给共享 DDR 和处理器：即使阵列内部 weight-stationary 很
高效，HWC 特征在每层前后的转置、展开或重新打包仍会出现在应用延迟中。
可重构数据流进一步允许按层选择循环映射，但“映射可变”并不自动解决外部
字节流的连续性、尾通道压紧和跨 pass 反馈。本文采用的 HWC 物化不是通用
NoC 重构，而是把软件张量语义保留到窗口生成边界，在该边界之后才形成固定
18-lane 向量。这个选择缩小了硬件支持范围，却使 DMA 输入、host golden 与
板端输出能够使用同一线性地址定义。

FPGA 编译框架的优势是从网络图自动生成专用流水或参数化核，但完整检测网络
会带来两个实际问题。第一，若每层各生成一套算子，早层大空间与深层大通道
会竞争片上资源，往往需要折叠或降低频率；第二，若复用一个卷积核，折叠后的
pass 边界、参数装载与中间张量布局就成为主要调度问题。LASA 属于第二类，
但它没有把折叠仅视为“吞吐下降的复用因子”，而是显式设计物化重放、反馈
所有权和上下文交接，以控制折叠产生的系统空泡。对于不能完整映射的历史
变体，本文只在代表层上讨论机制，不把局部加速外推为完整网络延迟。

目标检测专用 FPGA 工作还存在结果边界不统一的问题。有的工作报告骨干或
卷积核 GOPS，有的包含检测头但不包含 decode/NMS，有的使用不同输入尺度、
batch 或剪枝网络。即便都写作 YOLOv3-tiny，是否保留双尺度分支、255 通道头
和标准后处理也会显著改变计算量与精度。因此外部比较必须同时列出网络版本、
输入、batch、后处理和延迟边界。本文把完整性和可验证性单列为能力维度，
并不据此声称跨器件的绝对领先；真正可直接比较的受控对象是本文内部相同
RTL 根、相同阵列、相同时钟和相同输入下的消融变体。

FEATHER 提供了布局与数据流协同的重要参照，其核心启示是数据重排不应被
视为编译器之外的免费步骤。LASA 与之关注层次不同：本文不追求在多种片上
数据流之间任意切换，而是在完整检测网络所需的 HWC 接口下，构造一次物化、
多次重放的窄语义，并把跨 pass 的身份、所有权和顺序提交纳入协议。这样的
设计降低了通用性，却便于在 FPGA 上以有限 URAM、四路 DMA 和固定 200 MHz
时序实现，也使每个外部字节都能回溯到唯一的 layer/tile/pixel/channel。

\begin{table*}[t]
\centering
\caption{代表性方向的能力比较。符号“部分”表示论文可能支持该能力，但未同时给出完整网络、板端逐节点整数一致性与跨 pass 调度证据。}
\label{tab:related}
\small
\resizebox{0.98\textwidth}{!}{%
\begin{tabular}{lccccc}
\toprule
方向/代表工作 & 层间映射 & 外部HWC协同 & 跨pass重放 & 完整检测DAG & 逐节点板端证据\\
\midrule
固定数据流 ASIC（Eyeriss/TPU） & 部分 & 否/内部格式 & 否 & 否 & 否 \\
可重构数据流（MAERI/Simba） & 是 & 部分 & 部分 & 否 & 否 \\
FPGA 生成框架（FINN/fpgaConvNet） & 是 & 部分 & 部分 & 部分 & 部分 \\
YOLO FPGA 加速器 & 部分 & 部分 & 部分 & 是/不一 & 少见 \\
FEATHER & 是 & 是 & 非本文语义 & 非重点 & 否 \\
本文 LASA & 是 & 是 & 是 & 是 & 是 \\
\bottomrule
\end{tabular}
}
\end{table*}

